\chapter{Conclusions}\label{ch:conclusions}

This thesis evaluated three linear-time recurrent backbones
(RWKV-6, Mamba-2, Linear Attention) and three mechanism-level
extensions (MSDC, DHO, CVD) in a matched cross-architecture
transfer matrix. The results support mechanism-prior alignment:
mechanisms help most when their target axis matches both an
architectural deficit and the structure exercised by the task.

\section{Summary of findings}\label{sec:conclusions:summary}

The main empirical finding is that transfer is structured rather
than uniform. MSDC follows the local-mixing deficit ordering
LA $>$ RWKV-6 $>$ Mamba-2 across backbones. DHO produces its
strongest gains on Linear Attention, where the vanilla transition
has neither attenuation nor phase, but is distribution-modulated
on Mamba-2: the LibriSpeech near-null becomes productive on
Common Voice. CVD transfers asymmetrically and exposes decay as a
prerequisite in the bidirectional Linear Attention setting:
LION-LIT~$\times$~CVD is the only harmful cell in the matrix,
while controlled LION-S restores productive transfer. On MQAR,
MSDC and CVD solve the axis-2 recall probe at every reported
length while the causal Transformer baseline fails at $T = 1024$;
DHO fails as predicted by axis mismatch.

The best LibriSpeech cell is the 14M Mamba-2 CVD~$\times$~MSDC
composition at CER $0.0618$, but this number is secondary to the
cross-matrix pattern supporting the criterion.

\section{Limitations}\label{sec:conclusions:limitations}

Six limitations bracket the results. First, all matrix cells are
single-seed, so per-cell standard errors are not estimated; the
across-matrix claim rests on the coherence of many aligned
comparisons rather than on formal per-cell significance tests.
Second, the evidence is stronger on the positive-transfer side
of the criterion than on the full converse / null side;
systematic null-effect validation is left to future multi-seed
work, and the criterion is reported as a predictive heuristic
supported primarily on its positive-transfer half. Third, the
experiments do not fully separate structural inductive bias from
generic optimisation benefit; generic parameter-matched controls
(for example, MSDC compared to a residual block of equivalent
parameter count) are left for future work. Fourth, the matrix
stops at 14M parameters, so the deficit-driven ordering is not
asserted at scales beyond 14M. Fifth, the acoustic probes are
English-only, with Common Voice testing within-language
distribution shift rather than multilingual transfer. Sixth, the
mechanism hyperparameters (the MSDC dilation set
$\{1, 2, 4, 8\}$, the DHO depth-graded $\theta$-clip schedule of
$\pi/8$ and $\pi/4$, the CVD chunk-local preconditioner
temperature $\tau_{h}$) are reported as design choices following
literature precedent rather than as ablated optima.

\section{Scope and future work}\label{sec:conclusions:future-work}

The first extension is to cross the diagonal-class boundary. The
present criterion is tested only on diagonal-class recurrent
backbones. Repeating the matrix on diagonal-plus-low-rank models
such as DeltaNet~\cite{schlag2021deltanet, yang2024deltanet},
DeltaProduct~\cite{siems2025deltaproduct}, or
RWKV-7~\cite{peng2025rwkv7} would test whether the same alignment
logic holds when mechanisms target axis~3, non-Abelian
state-tracking. DeltaProduct is particularly suited because its
$n$-step Householder product provides a tunable state-tracking
capacity within the diagonal-plus-low-rank class.

The second extension is a stricter falsification protocol. The
criterion as stated tests a mechanism's gain against a post-hoc
reading of what a task exercises along each axis. A stronger
protocol would pre-register, for each
(architecture, mechanism, dataset) triple, a numerical prediction
derived from dataset-level metrics that quantify the structural
axes the dataset exercises (for example, a phoneme-to-syllable
energy spectrum for axis~1 on speech, a key-similarity entropy
for axis~2 on synthetic recall). The thesis does not develop
such metrics; their development is the natural prerequisite for
turning the alignment criterion from a predictive heuristic into
a strictly falsifiable claim.

The third extension is cross-modality replication. The same
matrix can be ported to vision and language probes by
substituting modality-appropriate axis-1 and axis-2 tasks for
the speech probes; the cross-modality scaffold is given by the
five-paper local-mixing cluster of
Section~\ref{sec:related:cross-modality}.

Overall, the thesis argues that mechanism transfer in linear-time
recurrent models is not generic. It is structured by the match
between the mechanism, the residual deficit of the backbone, and
the task structure exposed by the data. The matrix demonstrates
this structure within the diagonal-class family and motivates
extending the same test across larger scales, stricter
falsification protocols, and non-diagonal recurrent
architectures.
